\chapter{把动力学辨识写成可以计算的线性代数}\label{ch:identification}
\section{辨识不是重新发明机械臂的几何}
参数辨识（parameter identification）是利用测得的输入与运动，估计模型中的未知参数。URDF 或 DH 参数告诉我们连杆怎样连接、关节轴在哪里；它们使我们能够计算几何和动力学结构，但不保证每个质量、质心、转子惯量、摩擦和电流力矩系数都准确。辨识的任务是让这些数值与实际系统一致。

先区分三个问题。运动学标定修正关节零位和几何尺寸，主要解释位置误差；惯性参数辨识修正质量、质心和惯量，主要解释重力与加速力矩；摩擦辨识解释随速度、方向、温度和负载变化的阻力。三者可能互相污染，但它们不是同一个未知量。

本章以惯性与线性摩擦辨识为主线。目标不是让训练数据的误差尽量接近零，而是让没有用于拟合的新轨迹也能得到可靠力矩预测，并明确哪些参数实际可辨识。

\section{为什么对运动非线性 却对惯性参数线性}
\subsection{把参数选对 才能得到线性模型}
考虑绕水平轴转动的单杆。关节角从水平起量，质量 $m$，质心距离 $c$，质心惯量 $I_C$，则
\[
 \tau=(I_C+mc^2)\ddot q+mg_0c\cos q.
\]
若把未知量选为 $(m,c,I_C)$，式中出现 $mc^2$ 和 $mc$，并不是对这三个未知量线性。改用
\[
 I_O=I_C+mc^2,\qquad h=mc,
\]
就得到
\begin{equation}
 \tau=\begin{bmatrix}\ddot q&g_0\cos q\end{bmatrix}
 \begin{bmatrix}I_O\\h\end{bmatrix}.
\end{equation}
$h$ 叫一阶质量矩，单位 $\mathrm{kg\,m}$。运动量在前面的矩阵里，未知常数在后面的向量里。这种写法称为线性参数化，并不意味着运动方程对 $q$ 变成线性的。

对于三维刚体，用标准参数
\begin{equation}
 \pi_i=(m_i,h_{ix},h_{iy},h_{iz},I_{O,xx},I_{O,xy},I_{O,xz},I_{O,yy},I_{O,yz},I_{O,zz})^{\mathsf T}
\end{equation}
可以把刚体逆动力学写成 $\tau=Y(q,\dot q,\ddot q)\pi$。这里 $I_O$ 关于连杆坐标原点，而不是关于质心；所有惯量分量在同一连杆坐标系表达。库函数的参数顺序可能不同，调用时必须按其定义排列。

其根本原因是：刚体动能和重力势能对质量分布的积分是线性的，而欧拉拉格朗日的求导与相减也是线性运算。把质量积分归并成上述十个常数，剩余系数就仅依赖于几何和运动。

\subsection{一行数据到底包含什么}
加入库仑摩擦和黏性摩擦，对单关节可写
\begin{equation}\label{eq:single_reg_di}
 \tau_k=\underbrace{\begin{bmatrix}\ddot q_k&g_0\cos q_k&\operatorname{sgn}(\dot q_k)&\dot q_k\end{bmatrix}}_{H_k}
 \underbrace{\begin{bmatrix}I_O\\h\\F_c\\F_v\end{bmatrix}}_{\beta}+\epsilon_k.
\end{equation}
下标 $k$ 表示第 $k$ 次采样。$H_k$ 是回归行向量，$\beta$ 是待估参数，$\epsilon_k$ 是测量误差和模型遗漏的总残差。$F_c$ 单位 $\mathrm{N\,m}$，$F_v$ 单位 $\mathrm{N\,m\,s/rad}$。

例如在 $q=0$、$\dot q=0.4$、$\ddot q=1.2$ 时，$H_k=(1.2,9.81,1,0.4)$。如果 $\beta=(0.1,0.25,0.2,0.05)^{\mathsf T}$，则预测力矩为
\[
 1.2\times0.1+9.81\times0.25+1\times0.2+0.4\times0.05=2.7925\ \mathrm{N\,m}.
\]
每个数分别来自加速、重力、库仑摩擦和黏性摩擦，可以逐项核对。

七关节在一个采样时刻提供七行方程；采 $N$ 个时刻后，把所有行竖着堆叠，得到
\begin{equation}
 y=H\beta+\epsilon,\qquad y\in\mathbb R^{7N},\quad H\in\mathbb R^{7N\times p}.
\end{equation}
$p$ 为选用的参数个数。数据行很多不代表信息一定充分：把同一状态重复一万次，只会更精确地测到同一个方向的信息。

\section{基参数是可辨识的组合 不是某几根杆的原始参数}
\subsection{一个无法被任何轨迹分开的例子}
前面的单杆力矩只依赖于 $I_O$ 和 $h$。取以下两组参数：
\[
 (m,c,I_C)=(1,0.25,0.0375),\qquad(2,0.125,0.06875).
\]
它们都给出 $h=0.25$、$I_O=0.1$，因此在该单关节模型下，对任何 $q,\dot q,\ddot q$ 产生相同力矩。即使没有噪声，也不能仅凭这类关节运动数据区分这两组物理参数。这称为结构不可辨识。

基参数（base parameters）就是足以决定输入输出关系的一组独立参数组合。它们可用于力矩预测，但不一定能直接填写到某根连杆的质量和惯量字段里。

\subsection{用一个三列矩阵看清楚秩亏}
考虑
\[
 H=\begin{bmatrix}1&2&0\\2&4&1\\3&6&1\end{bmatrix}.
\]
第二列恒等于第一列的两倍。因此
\[
 H\begin{bmatrix}\beta_1\\\beta_2\\\beta_3\end{bmatrix}
 =\begin{bmatrix}1&0\\2&1\\3&1\end{bmatrix}
 \begin{bmatrix}\beta_1+2\beta_2\\\beta_3\end{bmatrix}.
\]
只能辨识 $\beta_{b1}=\beta_1+2\beta_2$ 和 $\beta_{b2}=\beta_3$。例如 $(1,2,3)$ 与 $(3,1,3)$ 的输出完全相同。矩阵有三列但秩只有二，直接计算 $(H^{\mathsf T}H)^{-1}$ 会失败。

\subsection{数值寻找基参数的方法}
在充分多样的无噪声模型状态上构造堆叠回归矩阵，用带列主元的 QR 分解选择独立列。若置换后的矩阵为 $[H_1\ H_2]$，其中 $H_2=H_1K$，则
\[
 H_1\pi_1+H_2\pi_2=H_1(\pi_1+K\pi_2)=H_1\pi_b.
\]
这里 $K$ 由线性方程求出；$\pi_b=\pi_1+K\pi_2$ 是基参数。SVD 同样可给出可辨识子空间，只是参数组合通常更稠密。

应同时检查：状态样本是否覆盖足够构型、各列是否合理缩放、阈值如何相对于最大奇异值设置。一个数值很小的奇异值可能表示结构依赖，也可能仅表示这次选取的状态没有激励到该方向。

\begin{pitfallbox}[七轴不自动等于四十三个惯性基参数]
基参数个数取决于关节类型、几何、基座约束、重力方向、是否含转子惯量以及参数化方式。某个七轴模型得到 $43$ 列，只说明那个模型和那套约定的结果；不能由“七轴”直接推断。每关节增加两列摩擦也不保证样本矩阵的秩恰好增加二：若速度只取一个非零值，$\operatorname{sgn}(\dot q)$ 与 $\dot q$ 列就成比例。
\end{pitfallbox}

\section{普通最小二乘从目标函数推到数字}
普通最小二乘 OLS 的目标是使平方残差之和最小：
\[
 J(\beta)=\|y-H\beta\|_2^2=(y-H\beta)^{\mathsf T}(y-H\beta).
\]
展开：
\[
 J=y^{\mathsf T}y-2\beta^{\mathsf T}H^{\mathsf T}y+\beta^{\mathsf T}H^{\mathsf T}H\beta.
\]
对 $\beta$ 求梯度并令其为零：
\[
 \nabla J=-2H^{\mathsf T}y+2H^{\mathsf T}H\hat\beta=0.
\]
因此正规方程是 $H^{\mathsf T}H\hat\beta=H^{\mathsf T}y$。仅当 $H$ 列满秩时，才能写出唯一解
\begin{equation}
 \hat\beta=(H^{\mathsf T}H)^{-1}H^{\mathsf T}y.
\end{equation}
这条式子解释数学，不推荐照着显式求逆实现。QR 或 SVD 可避免先构造病态程度被平方的矩阵。

\begin{examplebox}[算例 8A 四条数据完整求惯量与黏性系数]
假设已扣除重力与库仑摩擦，模型为 $\tau=J_r\ddot q+b\dot q$。数据矩阵与力矩为
\[
 H=\begin{bmatrix}1&1\\2&1\\1&2\\2&2\end{bmatrix},\qquad
 y=\begin{bmatrix}0.31\\0.39\\0.49\\0.62\end{bmatrix}.
\]
第一列是 $\mathrm{rad/s^2}$，第二列是 $\mathrm{rad/s}$，输出为 $\mathrm{N\,m}$。逐项相乘相加，得
\[
 H^{\mathsf T}H=\begin{bmatrix}10&9\\9&10\end{bmatrix},\qquad
 H^{\mathsf T}y=\begin{bmatrix}2.82\\2.92\end{bmatrix}.
\]
两阶逆矩阵的分母为 $10\times10-9\times9=19$，所以
\[
 \hat\beta=\frac1{19}\begin{bmatrix}10&-9\\-9&10\end{bmatrix}
 \begin{bmatrix}2.82\\2.92\end{bmatrix}
 =\begin{bmatrix}0.10105263\\0.20105263\end{bmatrix}.
\]
即 $\hat J_r=0.10105263\ \mathrm{kg\,m^2}$，$\hat b=0.20105263\ \mathrm{N\,m\,s/rad}$。再代回每条数据：
\[
 \hat y=(0.3021053,0.4031579,0.5031579,0.6042105)^{\mathsf T}.
\]
残差为 $r=y-\hat y=(0.0078947,-0.0131579,-0.0131579,0.0157895)^{\mathsf T}$。计算应满足 $H^{\mathsf T}r\approx0$：最优残差与回归矩阵列空间正交，但每个样本残差并不为零。
\end{examplebox}

\section{加权最小二乘的权重是噪声精度}
若各观测噪声方差不同，直接相加平方残差就会把不同可靠程度的数据当成一样可靠。设 $\epsilon$ 的协方差为 $\Sigma$，加权最小二乘 WLS 最小化
\begin{equation}
 J_W=(y-H\beta)^{\mathsf T}\Sigma^{-1}(y-H\beta).
\end{equation}
若噪声独立，$\Sigma=\operatorname{diag}(\sigma_1^2,\ldots,\sigma_m^2)$，选择
\[
 W=\operatorname{diag}(1/\sigma_1,\ldots,1/\sigma_m),
\]
则目标就是 $\|Wy-WH\beta\|^2$。也就是说，每行力矩和该行所有回归系数一起除以该行的噪声标准差。这里的 $W$ 是白化矩阵；目标函数中的权重是 $W^{\mathsf T}W=\Sigma^{-1}$。把方差与标准差混用，会产生错误的加权强度。

若噪声存在关节间或时间上的相关性，需要能满足 $W^{\mathsf T}W=\Sigma^{-1}$ 的一般白化变换，不再是简单逐行缩放。滤波后的残差通常存在时间相关性，不能自动把每个采样点当作独立观测。

\begin{examplebox}[算例 8B 两个测量值的最优加权平均]
两次测量同一个常数力矩，得到 $y_1=2$、$y_2=4\ \mathrm{N\,m}$；标准差分别为 $1$ 和 $2\ \mathrm{N\,m}$。目标函数
\[
 J=(2-\beta)^2+\frac14(4-\beta)^2.
\]
求导为 $2(\beta-2)+\tfrac12(\beta-4)=0$，得到
\[
 \hat\beta=\frac{2/1^2+4/2^2}{1/1^2+1/2^2}=2.4\ \mathrm{N\,m}.
\]
简单平均是 $3$；WLS 更相信标准差较小的第一个测量。权重来自噪声，不应仅因为某关节额定力矩大就人为降低它的重要性；若人为强调腕部精度，那是另一个设计目标，应与统计白化分开记录。
\end{examplebox}

\section{条件数 量纲与激励强度要一起看}
对列满秩矩阵，二范数条件数为 $\kappa_2(H)=\sigma_{\max}(H)/\sigma_{\min}(H)$。奇异值衡量不同参数方向在输出中的放大程度。最小奇异值很小时，某个参数组合对输出影响很弱，测量噪声容易掩盖它。

\subsection{为什么条件数会随单位改变}
例如 $H=\operatorname{diag}(1,0.01)$ 的条件数为 $100$。若第二个参数原来以米为单位，改用厘米，使相应回归列乘以 $0.01$，条件数就会进一步改变，物理实验本身却没有改变。带不同物理单位的列直接比较，数值可能主要反映单位选择。

选取固定参数尺度 $D=\operatorname{diag}(d_1,\ldots,d_p)$，定义 $\beta=D\vartheta$，则有
\begin{equation}
 H_s=WHD,\qquad y_s=Wy,\qquad y_s=H_s\vartheta.
\end{equation}
$\vartheta$ 是无量纲或相近量级的参数。计算条件数、比较轨迹时，$W$ 与 $D$ 应固定，否则每条轨迹各自重新归一化会掩盖绝对信息量的变化。

\subsection{条件数小并不等于信息多}
矩阵 $H_1=I_2$ 和 $H_2=0.001I_2$ 都有条件数一，但在相同输出噪声下，$H_2$ 的参数标准差是 $H_1$ 的一千倍。前者和后者“各方向一样均衡”，却不一样强。因此应联合观察最小奇异值、信息矩阵、参数协方差和实际力矩信号幅值。

当 $H$ 列满秩时，$H^{\mathsf T}H$ 的特征值是奇异值平方，所以
\[
 \kappa_2(H^{\mathsf T}H)=\kappa_2(H)^2.
\]
例如 $\kappa_2(H)=40$ 时，$\kappa_2(H^{\mathsf T}H)=1600$。若比较对象经过不同白化、缩放或剔除了不同列，这个数字就没有直接可比性。秩亏时最小奇异值为零，通常把条件数视为无穷大。

\begin{figure}[htbp]\centering
\includegraphics[width=0.91\linewidth]{figures/di_information_ellipses.pdf}
\caption{不同回归矩阵对应的参数不确定性椭圆。均衡程度由条件数描述，绝对大小还取决于总体激励强度。图中为合成二维线性模型。}
\end{figure}

\section{参数不确定性怎样计算 怎样正确解释}
在 $H$ 无误差、列满秩、噪声独立同方差且零均值的条件下，OLS 估计协方差为
\[
 \operatorname{Cov}(\hat\beta)=\sigma^2(H^{\mathsf T}H)^{-1}.
\]
如果 $\sigma$ 未知，可以用 $\hat\sigma^2=\|r\|^2/(m-p)$ 估计，$m$ 是观测行数，$p$ 是估计参数数目。分母扣掉 $p$，因为拟合参数消耗了自由度。模型偏差、滤波相关性或回归矩阵噪声存在时，这个公式不再给出完整不确定性。

算例 8A 中 $m=4,p=2$，残差平方和为 $0.0006578947$，所以
\[
 \hat\sigma=\sqrt{0.0006578947/2}=0.0181369\ \mathrm{N\,m}.
\]
参数协方差估计约为
\[
 \begin{bmatrix}0.00017313&-0.00015582\\-0.00015582&0.00017313\end{bmatrix}.
\]
两个参数标准差都约为 $0.013158$，但相对标准差不同，因为参数本身分别约为 $0.101$ 和 $0.201$。较大的负相关说明：惯量估高一些时，黏性系数估低一些，仍可部分保持当前数据的力矩预测。

新的运动状态给出回归行 $h_*$，参数不确定性引起的预测方差为
\begin{equation}
 \sigma_{\hat\tau,*}^2=h_*\operatorname{Cov}(\hat\beta)h_*^{\mathsf T}.
\end{equation}
如果要预测一次新的有噪声测量，还需加该测量噪声方差。只看每个参数的独立标准差而忽略参数相关性，会错误估计总预测误差。

\section{回归矩阵也有噪声时 最小二乘为何会偏}
前面的统计推导把 $H$ 当成准确已知。机器人数据中，$H$ 由位置、速度和加速度构造，它本身也受到编码器噪声、求导和时间错位影响。这叫变量含误差问题；增加样本数量可以降低随机波动，却不一定消除由此产生的系统偏差。

用一个标量模型看清楚原因。设真实关系为 $y=\beta x$，真实 $x$ 零均值、方差为 $s_x^2$，但实际记录的是 $\tilde x=x+\nu$；$\nu$ 为独立零均值误差，方差 $s_\nu^2$。用记录的 $\tilde x$ 做无截距最小二乘，很多样本时估计趋近于
\begin{equation}
 \hat\beta\approx\frac{\mathbb E[\tilde xy]}{\mathbb E[\tilde x^2]}
 =\frac{\beta\mathbb E[x^2]+\beta\mathbb E[\nu x]}{\mathbb E[x^2]+2\mathbb E[x\nu]+\mathbb E[\nu^2]}
 =\beta\frac{s_x^2}{s_x^2+s_\nu^2}.
\end{equation}
独立性使交叉项为零，但分母保留了额外噪声方差，导致估计向零衰减。这是一个明确的偏差，不是数值求逆不精确；换成更精确的线性代数求解器也不会消失。

\begin{examplebox}[算例 8F 加速度噪声可能把惯量估低]
假设只有惯性力矩 $y=J_r\ddot q$，真实惯量 $J_r=0.1$，真实加速度方差为 $1$，加速度估计误差方差为 $0.25$，且满足上述独立模型。则估计极限为
\[
 \hat J_r\approx0.1\frac1{1+0.25}=0.08\ \mathrm{kg\,m^2}.
\]
即使无限增加同类数据，仍可能低估 $20\%$。真实闭环系统的误差相关结构更复杂，偏差不一定总向零，但这个算例足以说明“数据很多”不能替代测量链质量。
\end{examplebox}

闭环采集还增加另一个相关性来源：测量噪声影响控制器输出，控制器输出又影响未来运动。因此不能简单认为力矩与位置噪声彼此独立。工具变量法尝试寻找一个与真实运动相关、但不与测量误差相关的替代信息；若工具本身来自同一个带噪信号，又没有验证独立性，就未必解决问题。

工程中应先减少可解释的误差：使用实际时间戳对齐，避免饱和样本，拟合实际周期运动，核对电流到力矩的比例，再评估是否需要更复杂估计方法。比较两个方法时，除了平均参数值，还要比较重复采集的波动、新轨迹预测和物理一致性，不能凭训练残差更小就认定偏差更小。

\section{物理一致性约束为什么比惯量正定更强}
\subsection{正惯量不一定来自真实质量分布}
惯量主值 $I_1,I_2,I_3$ 必须满足三角不等式：$I_1+I_2\ge I_3$，另外两种排列同样成立。例如 $\operatorname{diag}(1,1,3)$ 虽然正定，却不能作为任何非负质量分布的质心惯量，因为 $1+1<3$。

更统一的表达是伪惯量矩阵
\begin{equation}\label{eq:pseudoinertia_di}
 P_i=\begin{bmatrix}\Sigma_i&h_i\\h_i^{\mathsf T}&m_i\end{bmatrix},\qquad
 \Sigma_i=\tfrac12\operatorname{tr}(I_{O,i})I_3-I_{O,i}.
\end{equation}
这里 $\Sigma_i$ 是二阶质量矩，不是前面噪声协方差；本小节通过下标和定义区分两者。因为
\[
 P_i=\int\begin{bmatrix}r\\1\end{bmatrix}\begin{bmatrix}r^{\mathsf T}&1\end{bmatrix}dm,
\]
对任何向量 $z$，$z^{\mathsf T}P_iz=\int(z^{\mathsf T}[r;1])^2dm\ge0$，所以它必须半正定，并且质量为正。

对于理想零厚度杆或平面薄片，半正定矩阵可以是奇异的；严格正定对应更强的非退化空间质量分布。数值优化常施加一个小的正定裕量，但这属于建模与数值选择，不应声称所有理想物体都必须严格正定。

\begin{examplebox}[算例 8C 平行轴变换后的物理一致性]
质量 $m=2\ \mathrm{kg}$ 的均匀实心球半径 $0.1\ \mathrm m$，故 $I_C=(2mr^2/5)I_3=0.008I_3$。令其质心 $c=(0.2,0,0)^{\mathsf T}\ \mathrm m$，则
\[
 h=(0.4,0,0)^{\mathsf T},\qquad I_O=\operatorname{diag}(0.008,0.088,0.088).
\]
由式\eqref{eq:pseudoinertia_di}得 $\Sigma_i=\operatorname{diag}(0.084,0.004,0.004)$。消去右下角的质量块 $m$ 后，剩余矩阵称为 Schur 补；当 $m>0$ 时，其正定性等价于整个伪惯量的正定性。这里得到
\[
 \Sigma_i-\frac{hh^{\mathsf T}}m
 =\operatorname{diag}(0.084-0.16/2,0.004,0.004)=0.004I_3>0.
\]
因此该伪惯量严格正定。这里的减法把原点二阶矩还原成关于质心的二阶矩，物理意义十分直接。
\end{examplebox}

\subsection{为什么约束通常回到标准参数空间}
若最终要生成可用的仿真模型，应对每根连杆的标准参数施加约束。例如
\begin{equation}
 \min_{\pi}\ \|W(y-H_{\mathrm{full}}\pi)\|^2
 +\lambda\|D_0^{-1}(\pi-\pi_0)\|^2,\qquad P_i(\pi_i)\succeq0,\quad m_i>0.
\end{equation}
$\pi_0$ 是 CAD 先验，$D_0$ 用于量纲与先验尺度，$\lambda$ 控制对先验的信任。数据没有决定的参数方向，由先验和物理约束补充，不应宣称它们是被实验唯一测出来的。

质量为正、惯量符合三角不等式也不保证质心位于实际连杆包络内。若知道质量分布所在的几何范围，还需附加相应约束。检查若干姿态下 $M(q)>0$ 是有用的实现检查，却不能代替每根刚体参数的物理一致性论证。

\section{递推辨识和其他方法各自解决什么}
\subsection{递推最小二乘的一次完整更新}
当数据逐条到来，可用递推最小二乘 RLS 更新参数：
\begin{align}
 K_k&=\frac{P_{k-1}\varphi_k}{\lambda+\varphi_k^{\mathsf T}P_{k-1}\varphi_k},\\
 \hat\beta_k&=\hat\beta_{k-1}+K_k(y_k-\varphi_k^{\mathsf T}\hat\beta_{k-1}),\\
 P_k&=\lambda^{-1}(P_{k-1}-K_k\varphi_k^{\mathsf T}P_{k-1}).
\end{align}
$\varphi_k$ 是当前回归列向量，$\lambda\in(0,1]$ 是遗忘因子，$P$ 是信息矩阵的逆或其尺度版本。不能不加噪声尺度说明就总把 $P$ 当成真实参数协方差。

\begin{examplebox}[算例 8D 一次 RLS 更新]
只估一个参数，初值 $\hat\beta_0=2$，$P_0=1$，取 $\lambda=1$。新数据 $\varphi_1=2$、$y_1=6$。预测值为 $2\times2=4$，残差为 $2$；增益
\[
 K_1=\frac{1\times2}{1+2^2\times1}=0.4.
\]
所以 $\hat\beta_1=2+0.4\times2=2.8$，$P_1=1-0.4\times2=0.2$。它没有立刻变成 $6/2=3$，因为初始有限 $P_0$ 相当于保留先验。$\lambda=1$ 与批量最小二乘的等价，需要包含相同的初始先验约定。
\end{examplebox}

若 $\lambda<1$，旧数据逐渐被遗忘，可以跟踪慢变化参数，但长时间没有激励时，不确定性可能增长。不能只因为每周期都有新数据就每周期无条件更新；需要监测信息量、参数边界和应用效果。

\subsection{方法选择应对应具体误差来源}
\begin{longtable}{p{0.18\linewidth}p{0.34\linewidth}p{0.39\linewidth}}
\toprule 方法&主要用途&必须核对的假设\\\midrule\endhead
OLS 与 WLS&固定结构的线性参数估计&回归矩阵近似无误差，残差统计与权重匹配\\
正则化估计&激励不足时结合 CAD 等先验&先验会影响数据未确定的方向，需说明尺度\\
总体最小二乘 \mbox{TLS}&考虑回归矩阵与输出同时有误差&需要适合的噪声结构和量纲处理，普通 TLS 并非自动优于 OLS\\
工具变量法 IV&减少闭环与回归量噪声造成的偏差&工具变量要与真实回归量相关，并与误差不相关\\
最大似然 ML&利用明确的噪声概率模型&高斯输出噪声与已知协方差下可退化为 WLS\\
输出误差方法&通过仿真轨迹拟合位置或速度&通常是非线性优化，依赖初值、控制器模型与时间同步\\
能量或积分模型&减少对显式加速度的依赖&积分边界、摩擦功和可辨识参数集合会改变\\
残差学习&拟合已解释物理之外的重复误差&需要独立验证、输入范围约束，不保证范围外外推\\\bottomrule
\end{longtable}

\section{训练误差 验证误差与三种常见指标}
残差定义为 $e_k=\tau_k-\hat\tau_k$。RMSE 是均方根误差：
\[
 \mathrm{RMSE}=\sqrt{\frac1N\sum_{k=1}^N e_k^2}.
\]
它与力矩单位相同。所谓 RRMSE 没有唯一的跨文献定义，本书固定采用
\[
 \mathrm{RRMSE}_{\mathrm{rms}}=\frac{\sqrt{N^{-1}\sum e_k^2}}{\sqrt{N^{-1}\sum\tau_k^2}}\times100\%.
\]
如果别人用信号峰峰值、均值或去均值后的标准差作为分母，数值就不应直接比较。分母接近零时，相对误差可能巨大，应同时报告绝对 RMSE。

决定系数定义为
\[
 R^2=1-\frac{\sum e_k^2}{\sum(\tau_k-\bar\tau)^2},\qquad
 \bar\tau=N^{-1}\sum\tau_k.
\]
当分母非零时，$R^2=1$ 表示完美预测，$R^2=0$ 表示平方误差等同于始终预测这组数据的均值，$R^2<0$ 表示更差。验证集上的 $R^2$ 可以为负；常量信号的分母为零时没有通常意义下的该指标。

\begin{examplebox}[算例 8E 把三个指标都算一遍]
测量力矩 $\tau=(1,2,3,4)$，预测 $(1.1,1.8,3.2,3.9)$。误差为 $(-0.1,0.2,-0.2,0.1)$，平方和 $0.1$，因此
\[
 \mathrm{RMSE}=\sqrt{0.1/4}=0.158114\ \mathrm{N\,m}.
\]
测量力矩 RMS 是 $\sqrt{(1+4+9+16)/4}=\sqrt{7.5}=2.73861$，所以本书定义的 RRMSE 为 $5.7735\%$。均值为 $2.5$，去均值平方和为 $2.25+0.25+0.25+2.25=5$，所以 $R^2=1-0.1/5=0.98$。
\end{examplebox}

\begin{figure}[htbp]\centering
\includegraphics[width=0.91\linewidth]{figures/di_regression_validation.pdf}
\caption{合成验证轨迹中的模型预测与残差。残差与速度方向或加速度相关时，说明误差具有结构，不能仅用一个总体 RMSE 掩盖。}
\end{figure}

参数拟合、超参数选择和最终验证应使用可区分的数据。把同一周期随机拆成训练与验证点，可能因为强时间相关而高估泛化能力；完整独立轨迹、不同速度和姿态区域更能检验模型。误差是否可接受，应由最终任务决定：重力补偿看静态与低速力矩残差，快速跟踪看动态残差，外力观测还需考虑残差如何映射成假接触力。

\section{计算练习与分步解答}
\subsection*{练习 8.1 无法分开的摩擦参数}
只有恒速 $\dot q=0.2\ \mathrm{rad/s}$ 的数据，模型为 $F_c\operatorname{sgn}(\dot q)+F_v\dot q$。能否分别求 $F_c,F_v$？若测得摩擦为 $0.5\ \mathrm{N\,m}$，列举两组解。

\textbf{解答。}只能得到 $F_c+0.2F_v=0.5$，两列成比例，秩为一。例如 $(F_c,F_v)=(0.4,0.5)$ 或 $(0.3,1)$ 都成立。加入另一个不同非零速度，并确保其余条件一致，才可能把截距和斜率分开。

\subsection*{练习 8.2 白化必须同时处理输出和回归矩阵}
测量方程 $2\beta=6+\epsilon$，噪声标准差 $0.5$。写出白化后的方程。

\textbf{解答。}两边同除 $0.5$，得 $4\beta=12+\epsilon/0.5$，新噪声标准差为一。若只将输出变成 $12$ 而保留回归系数 $2$，就会错误估出 $\beta=6$，而原方程的解是 $3$。

\subsection*{练习 8.3 条件数与绝对误差}
$H=\operatorname{diag}(2,0.1)$，输出噪声为 $(0,0.01)^{\mathsf T}$。求条件数和参数误差。

\textbf{解答。}奇异值为 $2,0.1$，故条件数 $20$。参数误差 $H^{-1}\epsilon=(0,0.1)^{\mathsf T}$。若把整个 $H$ 乘十，条件数仍为 $20$，相同输出噪声导致的参数误差却变为 $(0,0.01)^{\mathsf T}$。这直接说明降低条件数与提高信号强度不是同一件事。

\subsection*{练习 8.4 主惯量检查}
质量为正，质心惯量主值分别为 $(0.01,0.02,0.04)\ \mathrm{kg\,m^2}$。它能否代表一个真实非负质量分布？

\textbf{解答。}虽然三个主惯量都为正，但 $0.01+0.02=0.03<0.04$，违反三角不等式，因此不能。不能只检查惯量矩阵可逆或正定就认定物理可行。

\section*{本章与原稿的对应}
本章展开 [A，第33--34页] 与 [B，第41--53页] 的线性参数化、基参数、最小二乘、协方差和辨识方法。原稿中针对某类机械臂的参数数量、仿真结果及经验阈值不作为普遍结论沿用；本章提供可逐项复算的小型数据集，并明确统计与物理假设。
